\documentclass{beamer}

\usepackage[T2A]{fontenc}
\usepackage[utf8]{inputenc}
\usepackage[english,russian]{babel}
\input{settings.tex}

\newcommand{\myqrframe}{
	\begin{frame}
		\centerline{Слайды доступны на Github:}
		\bigskip
		\centerline{\mygit}
		\bigskip
		\myqr
	\end{frame}
}


\subtitle{Современные методы управления роботами}
\author{Сергей Савин}
\date{2026}

\title{Оптимизация методом суммы квадратов}
\centering



\begin{document}
\maketitle

\begin{frame}{Содержание}
	\begin{itemize}
		\item Выпуклые функции
		\item Выпуклые области
		\item Выпуклая оптимизация
		\item Свойства выпуклой оптимизации
		\item Сумма квадратов
		\item Полиномиальный базис
		\item Матрица Грама
		\item Оптимизация методом суммы квадратов
		\item Оптимизация методом суммы квадратов в мономиальном базисе
	\end{itemize}
\end{frame}


\begin{frame}
	\begin{center}
		{\Large Часть 1: Выпуклая оптимизация}
	\end{center}
\end{frame}


\begin{frame}{Выпуклые функции}
	\begin{flushleft}
		Выпуклая функция имеет форму «чаши»:
		
		\bigskip
		
		Любой отрезок прямой, соединяющий две точки на её графике, целиком
		лежит \textbf{выше} (или касается) самого графика.
		
		\bigskip
		
		Эквивалентно, любая касательная к функции лежит
		\textbf{ниже} (или касается) её графика.
		
		\vspace{-2.5em}
		
		\begin{center}
			\begin{minipage}[t]{.4\textwidth}
				\centering
				\begin{tikzpicture}[x=0.75pt,y=0.75pt,yscale=-1,xscale=1]
					\draw [draw opacity=0, line width=1.0pt] (696.23,142.36) .. controls (681.43,167.71) and (649.5,185.54) .. (612.24,186.11) .. controls (573.51,186.7) and (539.97,168.5) .. (525.19,141.99) -- (611.08,111.08) -- cycle ;
					\draw (696.23,142.36) .. controls (681.43,167.71) and (649.5,185.54) .. (612.24,186.11) .. controls (573.51,186.7) and (539.97,168.5) .. (525.19,141.99) ;
					\draw [color={rgb, 255:red, 208; green, 2; blue, 27}, draw opacity=1, line width=1.75pt] (680,161.5) -- (560,173.5) ;
					\draw [color={rgb, 255:red, 208; green, 2; blue, 27}, draw opacity=1] (680,161.5) -- (560,173.5) ;
					%Shape: Circle [id:dp3160827101885906]
					\draw [fill={rgb, 255:red, 0; green, 0; blue, 0}, fill opacity=1] (554.5,173.5) .. controls (554.5,170.46) and (556.96,168) .. (560,168) .. controls (563.04,168) and (565.5,170.46) .. (565.5,173.5) .. controls (565.5,176.54) and (563.04,179) .. (560,179) .. controls (556.96,179) and (554.5,176.54) .. (554.5,173.5) -- cycle ;
					%Shape: Circle [id:dp16415395117171228]
					\draw [fill={rgb, 255:red, 0; green, 0; blue, 0}, fill opacity=1] (674.5,161.5) .. controls (674.5,158.46) and (676.96,156) .. (680,156) .. controls (683.04,156) and (685.5,158.46) .. (685.5,161.5) .. controls (685.5,164.54) and (683.04,167) .. (680,167) .. controls (676.96,167) and (674.5,164.54) .. (674.5,161.5) -- cycle ;
					\draw (560,185) node [anchor=north west, inner sep=0.75pt, align=left] {$x_1$};
					\draw (675,170) node [anchor=north west, inner sep=0.75pt, align=left] {$x_2$};
				\end{tikzpicture}
				\captionof{figure}{Выпуклая функция}
			\end{minipage}%
			\begin{minipage}[t]{.7\textwidth}
				\centering
				\begin{tikzpicture}[x=0.75pt,y=0.75pt,yscale=-1,xscale=1]
					%Straight Lines [id:da6875192047258083]
					\draw [color={rgb, 255:red, 208; green, 2; blue, 27}, draw opacity=1, line width=1.75pt] (564.6,165.5) -- (444.6,177.5) ;
					%Shape: Circle [id:dp14324656781298462]
					\draw [fill={rgb, 255:red, 0; green, 0; blue, 0}, fill opacity=1] (439.1,177.5) .. controls (439.1,174.46) and (441.56,172) .. (444.6,172) .. controls (447.64,172) and (450.1,174.46) .. (450.1,177.5) .. controls (450.1,180.54) and (447.64,183) .. (444.6,183) .. controls (441.56,183) and (439.1,180.54) .. (439.1,177.5) -- cycle ;
					%Shape: Circle [id:dp2582069443071713]
					\draw [fill={rgb, 255:red, 0; green, 0; blue, 0}, fill opacity=1] (559.1,165.5) .. controls (559.1,162.46) and (561.56,160) .. (564.6,160) .. controls (567.64,160) and (570.1,162.46) .. (570.1,165.5) .. controls (570.1,168.54) and (567.64,171) .. (564.6,171) .. controls (561.56,171) and (559.1,168.54) .. (559.1,165.5) -- cycle ;
					%Curve Lines [id:da47039197462224447]
					\draw (498.2,186.8) .. controls (512.6,148) and (529.4,88) .. (547.4,128) ;
					%Curve Lines [id:da13389450894014865]
					\draw (399,161.2) .. controls (439,131.2) and (470.6,244.8) .. (498.2,186.8) ;
					%Curve Lines [id:da2313253314744761]
					\draw (547.4,128) .. controls (563.8,173.6) and (577.4,192) .. (617.4,162) ;
					\draw (440,190) node [anchor=north west, inner sep=0.75pt, align=left] {$x_1$};
					\draw (560,180) node [anchor=north west, inner sep=0.75pt, align=left] {$x_2$};
				\end{tikzpicture}
				\captionof{figure}{Невыпуклая функция}
			\end{minipage}
		\end{center}
		
	\end{flushleft}
\end{frame}


\begin{frame}{Выпуклые области}
	\begin{flushleft}
		
		Выпуклая область — это форма или пространство, в котором отрезок прямой,
		соединяющий любые две точки внутри области, целиком остаётся внутри неё.
		
		\bigskip
		
		\begin{figure}
			\centering
			\begin{minipage}{.3\textwidth}
				\centering
				
				\textit{  
					\begin{tikzpicture}[x=0.75pt,y=0.75pt,yscale=-1,xscale=1]
						\draw  [fill={rgb, 255:red, 255; green, 160; blue, 160 }  ,fill opacity=1 ] (148,55) -- (188,62.5) -- (211,107.5) -- (175,170.5) -- (94,158.5) -- (98,85) -- cycle ;
				\end{tikzpicture}}
				
				\captionof{figure}{Выпуклая область}
			\end{minipage}%
			\begin{minipage}{.7\textwidth}
				\centering
				
				\begin{tikzpicture}[x=0.75pt,y=0.75pt,yscale=-1,xscale=1]
					\draw  [fill={rgb, 255:red, 255; green, 160; blue, 160 }  ,fill opacity=1 ] (525,68.5) -- (525,112.5) -- (452,97.5) -- (435,120.5) -- (533,188.5) -- (488,204.5) -- (411,187.5) -- (404,72.5) -- cycle ;
					\draw [color={rgb, 255:red, 208; green, 2; blue, 27 }  ,draw opacity=1, line width=1.75pt ]   (501,88.5) -- (486.03,126.12) -- (464,181.5) ;
				\end{tikzpicture}
				
				\captionof{figure}{Невыпуклая область}
			\end{minipage}
		\end{figure}
		
	\end{flushleft}
\end{frame}


\begin{frame}{Выпуклая оптимизация}
	\begin{flushleft}
		Задача вида
		\[
		\begin{aligned}
			\underset{x}{\text{minimize}}\quad & f(x)\\
			\text{subject to}\quad & g_i(x) \leq 0,\\
			& h_j(x) = 0
		\end{aligned}
		\]
		называется \textbf{выпуклой}, если
		
		\begin{itemize}
			\item $f(x)$ — выпуклая функция,
			\item каждая $g_i(x)$ — выпуклая функция,
			\item каждая $h_j(x)$ — аффинная функция.
		\end{itemize}
		
		\bigskip
		
		Например:
		\[
		\begin{aligned}
			\underset{x}{\text{minimize}}\quad &(x+1)^2\\
			\text{subject to}\quad &x \geq 0.
		\end{aligned}
		\]
		
	\end{flushleft}
\end{frame}


\begin{frame}{Свойства выпуклой оптимизации}
	\begin{flushleft}
		\begin{itemize}
			\item \textbf{Локальный минимум является глобальным минимумом}.
			
			\bigskip
			
			\item Выпуклые задачи часто могут решаться \textbf{эффективно}, даже при
			большом количестве переменных.
			
			\bigskip
			
			\item Существуют \textbf{надёжные алгоритмы} с сильными теоретическими
			гарантиями.
			
			\bigskip
			
			\item Существуют зрелые \textbf{решатели}: программные пакеты, предназначенные
			для решения определённых классов задач выпуклой оптимизации.
		\end{itemize}
		
	\end{flushleft}
\end{frame}


\begin{frame}
	\begin{center}
		{\Large Часть 2: Сумма квадратов}
	\end{center}
\end{frame}


\begin{frame}{Сумма квадратов}
	\begin{flushleft}
		Некоторые полиномы могут быть представлены как \emph{сумма квадратов} (SoS):
		\[
		p(x) = \sum_i g_i(x)^2,
		\]
		где $g_i(x)$ сами являются полиномами. Полином, представимый в виде суммы квадратов, всегда является
		\textbf{неотрицательным}:
		\[
		p(x) \geq 0 \qquad \text{для всех }x.
		\]
		
		
		\begin{itemize}
			\item Для произвольного полинома проверка его неотрицательности
			может быть вычислительно сложной (в общем случае NP-трудной).
			
			\item Проверка того, является ли полином суммой квадратов, может быть сведена к
			\textbf{задаче выпуклой оптимизации}.
			
			\item Это даёт практический способ нахождения
			\textbf{сертификатов неотрицательности}.
		\end{itemize}
		
	\end{flushleft}
\end{frame}


\begin{frame}{Полиномиальный базис}
	\begin{flushleft}
		
		Определим набор полиномов
		\[
		g_1(x),\;g_2(x),\;\ldots,\;g_N(x)
		\]
		и объединим их в вектор
		\[
		g(x)=
		\begin{bmatrix}
			g_1(x)\\
			\vdots\\
			g_N(x)
		\end{bmatrix}.
		\]
		
		\bigskip
		
		Используя этот вектор, можно строить полиномы из произведений
		$g_i(x)g_j(x)$ с весами $q_{ij}$:
		\[
		p(x)=\sum_{i=1}^N\sum_{j=1}^N
		q_{ij}g_i(x)g_j(x).
		\]
		
		\bigskip
		
		То же самое в матричной форме:
		\[
		p(x)=g(x)^\top Qg(x).
		\]
		
		
	\end{flushleft}
\end{frame}


\begin{frame}{Полиномиальный базис: пример}
	\begin{flushleft}
		
		Рассмотрим полиномиальный базис $g(x)$ и матрицу Грама $Q$:
		
		\[
		g(x)=
		\begin{bmatrix}
			1\\x
		\end{bmatrix},
		\qquad
		Q=
		\begin{bmatrix}
			2&1\\
			1&2
		\end{bmatrix}.
		\]
		
		Тогда мы можем построить полином:
		\[
		\begin{aligned}
			p(x)
			&=
			\begin{bmatrix}1&x\end{bmatrix}
			\begin{bmatrix}2&1\\1&2\end{bmatrix}
			\begin{bmatrix}1\\x\end{bmatrix}\\
			&=2+2x+2x^2.
		\end{aligned}
		\]
		
	\end{flushleft}
\end{frame}


\begin{frame}{Матрица Грама}
	\begin{flushleft}
		В выражении $p(x)=g(x)^\top Qg(x)$ матрица $Q$ называется
		\textbf{матрицей Грама}.
		
		\begin{theorem}
			Если $Q$ симметрична и положительно полуопределена, $Q \succeq 0$, то $p(x)=g(x)^\top Qg(x) \geq 0$ для всех $x$.
		\end{theorem}
		
		\textbf{Доказательство:} Для любого $x$ положим $v=g(x)$. Поскольку $Q\succeq0$, имеем $v^\top Qv\geq0$. Следовательно, $p(x)=g(x)^\top Qg(x)\geq0.$ \qed
		
		\bigskip
		
		Если $Q\succ0$ и $g(x)\neq0$ для всех $x$, то $p(x)>0$ для всех $x$.
		
	\end{flushleft}
\end{frame}


\begin{frame}{Матрица Грама: пример}
	\begin{flushleft}
		
		Для
		\[
		g(x)=
		\begin{bmatrix}
			1\\x
		\end{bmatrix},
		\qquad
		Q=
		\begin{bmatrix}
			2&1\\
			1&2
		\end{bmatrix},
		\]
		собственные значения $Q$ равны
		\[
		\lambda_1=2+1=3>0,
		\qquad
		\lambda_2=2-1=1>0.\footnotemark
		\]
		%
		Все собственные значения положительны, поэтому $Q\succ0$ (матрица положительно определена).
		
		
		\bigskip
		
		Поскольку $g(x)\neq0$ для любого $x$, то $p(x)=g(x)^\top Qg(x)>0$ для всех $x$. Действительно:
		\[
		p(x)=2+2x+2x^2>0.
		\]
		
		
		\footnotetext{\vspace{1.0em}
			Для матриц вида
			$\begin{bmatrix}a&b\\b&a\end{bmatrix}$,
			собственные значения равны $a+b$ и $a-b$.}
		
	\end{flushleft}
\end{frame}


\begin{frame}{Пример: проверка положительной определённости}
	\begin{flushleft}
		Рассмотрим $p(x)=3x_1^2+4x_1x_2+3x_2^2+5$. Используя мономиальный базис:
		\[
		v(x)=
		\begin{bmatrix}
			x_1 & x_2 & 1
		\end{bmatrix}^\top,
		\]
		полином можно записать в виде:
		\[
		p(x)=v(x)^\top Qv(x), \quad 
		Q=
		\begin{bmatrix}
			3&2&0\\
			2&3&0\\
			0&0&5
		\end{bmatrix}.
		\]
		
		Верхняя левая подматрица имеет собственные значения
		\[
		3+2=5>0,
		\qquad
		3-2=1>0.
		\]
		
		Собственные значения блочно-диагональной матрицы являются собственными значениями
		её блоков, $\lambda(Q)=\{5,1,5\}$, поэтому $Q \succ 0$. Следовательно:
		\[
		p(x)=v(x)^\top Qv(x)>0, \quad \forall x.
		\]
		
	\end{flushleft}
\end{frame}


\begin{frame}{Оптимизация методом суммы квадратов}
	\begin{flushleft}
		Оптимизация методом суммы квадратов (SoS) представляет собой класс задач выпуклой оптимизации,
		которые могут включать ограничения вида
		\[
		p(x) \in \mathrm{SoS}.
		\]
		
		\medskip
		
		Это означает, что мы ищем матрицу $Q$, такую что
		\[
		Q\succeq0
		\]
		и
		\[
		p(x)=z(x)^\top Qz(x),
		\]
		где $z(x)$ — фиксированный полиномиальный базис.
		
		\medskip
		
		Также можно потребовать
		\[
		Q\succ0
		\]
		если нам нужен строго положительный полином.
		
	\end{flushleft}
	
	\footnotetext{\vspace{1.0em}
		$Q \succ 0$ часто реализуется как
		$Q\succeq\epsilon I$ для малого $\epsilon>0$.}
\end{frame}


\begin{frame}{Оптимизация методом суммы квадратов в мономиальном базисе}
	\begin{flushleft}
		Распространённый выбор состоит в построении $z(x)$ из мономов до
		некоторой степени $N$:
		\[
		z(x)=
		\begin{bmatrix}
			1 & x_1 & x_2 & x_1x_2 & x_1^2 & x_2^2 & \cdots
		\end{bmatrix}^{\top}.
		\]
		
		\bigskip
		
		Если $z(x)$ содержит мономы до степени $N$, то
		\[
		z(x)^\top Qz(x)
		\]
		может содержать мономы до степени $2N$.
		
		\bigskip
		
		Например, полином степени $2N$ можно записать как
		\[
		p(x)=c_0+c_1x_1+c_2x_2+c_3x_1x_2
		+c_4x_1^2+\cdots.
		\]
		
		\bigskip
		
		Раскрываем выражение $z(x)^\top Qz(x)$ и \textbf{сопоставляем коэффициенты} соответствующих
		мономов.
		
		\bigskip
		
		Это даёт линейные ограничения равенства на элементы $Q$.
		
	\end{flushleft}
\end{frame}


\begin{frame}{Мономиальный базис — сопоставление коэффициентов}
	\begin{flushleft}
		
		Для заданного мономиального базиса $z(x)$ степени $N$ и монома $p_i(x)$ степени $\leq 2N$ в общем случае существует несколько способов воспроизвести этот моном через равенство $z(x)^\top Qz(x) = p_i(x)$.
		
		\bigskip
		
		Пример. $p_i(x) = x_1^4 x_2^2 x_3$ (моном степени 7), а $z(x)$ содержит все мономы вида $x_1^\alpha x_2^\beta x_3^\gamma$ степени не выше $4$. Тогда $p_i(x) $ можно представить как:
		
		\begin{align*}
			x_1^4 x_2^2 x_3 = 
			\textcolor{myblue}{x_1^4} \cdot \textcolor{mydarkgreen}{x_2^2 x_3} = 
			\textcolor{myblue}{x_1^2 x_2} \cdot \textcolor{mydarkgreen}{x_1^2 x_2 x_3} = 
			\textcolor{myblue}{x_1 x_2 x_3} \cdot \textcolor{mydarkgreen}{x_1^3 x_2} = \\ =
			\textcolor{myblue}{x_1^2 x_3} \cdot \textcolor{mydarkgreen}{x_1^2 x_2^2} = 
			\textcolor{myblue}{x_1^3} \cdot \textcolor{mydarkgreen}{x_1 x_2^2 x_3}
		\end{align*}
		
		В программе SoS коэффициент при $p_i(x) = x_1^4 x_2^2 x_3$ будет сопоставлен с суммой коэффициентов $q_{ij}$, соответствующих всем указанным выше произведениям.
		
		
	\end{flushleft}
\end{frame}


\begin{frame}{Ограничения в программе SoS: пример}
	\begin{flushleft}
		
		Рассмотрим полином $p(x)=a x_1^3+b x_1x_2^2+c x_2^4$. Выберем мономиальный базис:		
		\[
		z(x)=
		\begin{bmatrix}
			1 &
			x_1 &
			x_2 &
			x_1x_2 &
			x_1^2 &
			x_2^2
		\end{bmatrix}^\top
		\]
		%
		Затем
		\[
		\begin{bmatrix}
			1\\
			x_1\\
			x_2\\
			x_1x_2\\
			x_1^2\\
			x_2^2
		\end{bmatrix}^\top
		\begin{bmatrix}
			q_{11}&q_{12}&q_{13}&q_{14}&q_{15}&q_{16}\\
			q_{21}&q_{22}&q_{23}&q_{24}&q_{25}&\color{red}{q_{26}}\\
			q_{31}&q_{32}&q_{33}&\color{red}{q_{34}}&q_{35}&q_{36}\\
			q_{41}&q_{42}&\color{red}{q_{43}}&q_{44}&q_{45}&q_{46}\\
			q_{51}&q_{52}&q_{53}&q_{54}&q_{55}&q_{56}\\
			q_{61}&\color{red}{q_{62}}&q_{63}&q_{64}&q_{65}&q_{66}
		\end{bmatrix}
		\begin{bmatrix}
			1\\
			x_1\\
			x_2\\
			x_1x_2\\
			x_1^2\\
			x_2^2
		\end{bmatrix}
		=
		a x_1^3+\textcolor{red}{b}x_1 x_2^2 + c x_2^4
		\]
		
		Получаем следующие ограничения:
		
		\begin{align*}
			& q_{25} + q_{52} = a, & q_{26} + q_{62} + q_{34} + q_{43} = b, \\
			& q_{66} = c, & Q \succeq 0
		\end{align*}
		
	\end{flushleft}
\end{frame}


\begin{frame}{Literature}
	
	\begin{itemize}
		
		\item MIT OpenCourseWare -- Algebraic Techniques and Semidefinite Optimization:
		\bref{https://ocw.mit.edu/courses/6-972-algebraic-techniques-and-semidefinite-optimization-spring-2006/resources/lecture_10/}
		{Lecture 10: Sums of Squares and Semidefinite Programming.}
		
		\item Antonis Papachristodoulou and Stephen Prajna:
		\bref{https://users.ox.ac.uk/~engs0587/PapPACC05.pdf}
		{A Tutorial on Sum of Squares Techniques for Systems Analysis.}
		
		
	\end{itemize}
	
\end{frame}


\myqrframe


\begin{frame}{Приложение. Собственные значения симметричной матрицы $2 \times 2$}
	\begin{flushleft}
		Матрицы вида:
		\[
		\mathbf A = \begin{bmatrix} a & b \\ b & a \end{bmatrix}
		\]
		имеют собственные значения $\lambda_1 = a+b$ и $\lambda_2 = a-b$.
		
		\bigskip
		
		\textbf{Доказательство.} Непосредственное умножение матрицы на вектор $\mathbf A\mathbf{v} = \lambda \mathbf{v}$:
		\[
		\begin{aligned}
			\begin{bmatrix} a & b \\ b & a \end{bmatrix} \begin{bmatrix} 1 \\ 1 \end{bmatrix} &= \begin{bmatrix} a+b \\ a+b \end{bmatrix} = (a+b) \begin{bmatrix} 1 \\ 1 \end{bmatrix} \implies \lambda_1 = a+b\\[8pt]
			\begin{bmatrix} a & b \\ b & a \end{bmatrix} \begin{bmatrix} 1 \\ -1 \end{bmatrix} &= \begin{bmatrix} a-b \\ b-a \end{bmatrix} = (a-b) \begin{bmatrix} 1 \\ -1 \end{bmatrix} \implies \lambda_2 = a-b
		\end{aligned}
		\]
		
	\end{flushleft}
\end{frame}


\end{document}
